\section*{Methodology}

\subsection*{3.1 Study Population and Data Acquisition}

Seven patients with cerebrovascular pathology were prospectively enrolled and imaged using two MRI platforms: a clinical 3T Siemens system and an investigational 7T Siemens system. 4D Flow MRI enables simultaneous acquisition of three-directional velocity fields and is increasingly used for cerebrovascular hemodynamics assessment \cite{markl2012,dyverfeldt2015}. For each patient, paired 4D Flow MRI studies were acquired on both systems within a single session to minimize inter-session physiological variability. Each acquisition comprised four volumetric time series: a magnitude image and three directional phase-contrast velocity maps encoding the velocity components along the feet-head ($u$), left-right ($v$), and anterior-posterior ($w$) directions. The 3T acquisitions covered between 9 and 14 cardiac frames, while the 7T acquisitions spanned 10 to 15 frames, with a consistent velocity encoding of $\text{VENC} = 0.9$~m/s applied uniformly across all directional components and both field strengths \cite{bock2010}. All raw data were converted from DICOM to NIfTI format prior to further processing.

% FIGURE 1 — Acquisition and Pairing Overview:
% Schematic diagram illustrating the dual-field-strength acquisition protocol.
% For each patient, a representative axial slice from the 3T and 7T magnitude images
% should be shown side by side, highlighting differences in spatial resolution, SNR,
% and vascular conspicuity. A sagittal MIP of each acquisition below to convey 3D coverage.
\begin{figure}[h]
    \centering
    % Insert figure here
    \caption{Acquisition and Pairing Overview. Schematic diagram illustrating the dual-field-strength acquisition protocol.}
    \label{fig:acquisition}
\end{figure}

\subsection*{3.2 Motion Correction and Multi-Field-Strength Registration}

\subsubsection*{3.2.1 Intra-Scan Motion Correction}

Cardiac-gated 4D Flow MRI is intrinsically sensitive to bulk subject motion between cardiac phases \cite{stehning2020}. To correct for this, rigid-body transformations were estimated from the magnitude time series for each frame relative to a reference (first) frame, using normalized mutual information as the similarity criterion \cite{maes1997}. The resulting transformation matrices $\{\mathbf{T}_t^{\text{rigid}}\}_{t=1}^{T}$ were then applied consistently to the corresponding velocity phase images at each time point $t$:

\begin{equation}
    \mathbf{V}_t^{\text{corr}} = \mathbf{T}_t^{\text{rigid}} \circ \mathbf{V}_t^{\text{raw}}, \quad t = 1, \ldots, T
\end{equation}

where $\mathbf{V}_t^{\text{raw}} = (u_t, v_t, w_t)$ denotes the raw velocity volume at frame $t$. Following transformation, the velocity vectors were reoriented to maintain physical consistency in the scanner coordinate frame.

\subsubsection*{3.2.2 Inter-Scan Registration (7T $\rightarrow$ 3T)}

To align 7T volumes into the geometric space of the clinical 3T acquisitions, a multi-stage inter-scan registration was performed \cite{qu2020synth7t,cui2024synth7t}. A temporal maximum intensity projection (MIP) was first computed for each acquisition by retaining, at each voxel, the maximum magnitude value across all cardiac frames:

\begin{equation}
    \text{MIP}(\mathbf{x}) = \max_{t \in \{1,\ldots,T\}} M_t(\mathbf{x})
\end{equation}

where $M_t(\mathbf{x})$ denotes the magnitude volume at voxel $\mathbf{x}$ and time $t$. The MIP images provide a static, high-contrast angiographic representation suitable for robust registration. A two-stage registration strategy was employed: an initial rigid alignment $\mathbf{T}^{\text{rigid}}_{7T \to 3T}$, followed by an affine refinement $\mathbf{T}^{\text{affine}}_{7T \to 3T}$, both optimized using normalized mutual information. The composite transformation was then applied to the full 4D velocity volumes:

\begin{equation}
    \mathbf{V}_{7T \to 3T}^{\text{reg}}(\mathbf{x}, t) = \mathbf{T}^{\text{affine}}_{7T \to 3T} \circ \mathbf{V}_{7T}^{\text{corr}}(\mathbf{x}, t)
\end{equation}

This registration pipeline ensured voxel-level correspondence between paired 3T and 7T acquisitions, which is a prerequisite for supervised learning.

% FIGURE 2 — Registration Pipeline:
% Multi-panel figure: (a) temporal MIP of a 3T magnitude image; (b) temporal MIP of the
% corresponding 7T magnitude image before registration; (c) 7T MIP after rigid + affine
% registration into 3T space, with a checkerboard or difference overlay to demonstrate
% alignment quality. An additional panel showing the displacement field magnitude can be included.
\begin{figure}[h]
    \centering
    % Insert figure here
    \caption{Registration Pipeline. Multi-panel figure showing the temporal MIPs and registration results.}
    \label{fig:registration}
\end{figure}

\subsection*{3.3 Angiographic Volume Synthesis}

To produce a voxel-level vascular probability map suitable for vessel segmentation, a synthetic 3D angiographic representation was constructed by combining complementary contrast sources. First, a 4D speed image was derived by computing the voxel-wise Euclidean norm of the VENC-scaled velocity vector at each cardiac frame:

\begin{equation}
    S(\mathbf{x}, t) = \sqrt{
        \left(\frac{u(\mathbf{x},t)}{\text{VENC}}\right)^2 +
        \left(\frac{v(\mathbf{x},t)}{\text{VENC}}\right)^2 +
        \left(\frac{w(\mathbf{x},t)}{\text{VENC}}\right)^2
    }
\end{equation}

The temporal projection of the speed image $\bar{S}(\mathbf{x}) = \max_t S(\mathbf{x}, t)$ was then normalized to the unit interval. Likewise, the temporal MIP of the magnitude image $\overline{\text{MIP}}(\mathbf{x})$ was min-max normalized. The final angiographic representation $\mathcal{A}(\mathbf{x})$ was obtained as the element-wise product of these two normalized fields:

\begin{equation}
    \mathcal{A}(\mathbf{x}) = \hat{\text{MIP}}(\mathbf{x}) \cdot \hat{S}(\mathbf{x})
\end{equation}

where $\hat{\cdot}$ denotes min-max normalization. This product reinforces voxels that are simultaneously bright in magnitude and fast in flow velocity, yielding enhanced arterial contrast while suppressing static tissue and background noise.

% FIGURE 3 — Angiographic Volume Synthesis:
% Three-column panel showing axial and coronal views of: (a) normalized temporal MIP of the
% magnitude; (b) normalized temporal speed projection; (c) synthesized angiographic product
% A(x). A 3D volume rendering of A clearly delineating the Circle of Willis anatomy.
\begin{figure}[h]
    \centering
    % Insert figure here
    \caption{Angiographic Volume Synthesis. Three-column panel showing the normalized MIP, speed projection, and synthesized angiographic product $\mathcal{A}(\mathbf{x})$.}
    \label{fig:angiographic}
\end{figure}

\subsection*{3.4 Vascular Segmentation and Mask Generation}

\subsubsection*{3.4.1 Dual-Branch Segmentation Strategy}

Vessel masking was performed using a dual-branch strategy applied to the synthesized angiographic volume $\mathcal{A}(\mathbf{x})$, combining the complementary strengths of deep learning and classical image analysis.

\textbf{Deep Learning Branch.} A topology-aware nnU-Net ensemble from the TopCoW 2024 winning solution (CLAIM Team, Charité) \cite{yang2023topcow,aydin2024topcow} was applied to $\mathcal{A}(\mathbf{x})$ to produce a probabilistic segmentation $P^{\text{DL}}(\mathbf{x}) \in [0,1]$ of the Circle of Willis (CoW). This model was specifically trained to preserve topological consistency of the CoW vascular graph, producing anatomically coherent segmentations even in regions of low SNR.

\textbf{Classical Vesselness Branch.} In parallel, multi-scale Frangi and Sato vesselness filters \cite{frangi1998,sato1998} were computed over $\mathcal{A}(\mathbf{x})$. The Frangi filter exploits the eigenvalues $\lambda_1 \leq \lambda_2 \leq \lambda_3$ of the Hessian matrix $\mathbf{H}(\mathbf{x}, \sigma)$ at scale $\sigma$:

\begin{equation}
    V_\sigma(\mathbf{x}) = \begin{cases}
        0 & \text{if } \lambda_2 > 0 \text{ or } \lambda_3 > 0 \\[6pt]
        \left(1 - e^{-\mathcal{R}_A^2 / 2\alpha^2}\right)
        e^{-\mathcal{R}_B^2 / 2\beta^2}
        \left(1 - e^{-\mathcal{S}^2 / 2c^2}\right) & \text{otherwise}
    \end{cases}
\end{equation}

where $\mathcal{R}_A = |\lambda_2|/|\lambda_3|$, $\mathcal{R}_B = |\lambda_1|/\sqrt{|\lambda_2 \lambda_3|}$, and $\mathcal{S} = \|\mathbf{H}\|_F$ are shape and structure descriptors, and $\alpha, \beta, c$ are sensitivity parameters. The final vesselness response was obtained as the maximum across scales: $V(\mathbf{x}) = \max_\sigma V_\sigma(\mathbf{x})$. The resulting map was binarized via percentile-based thresholding and refined through morphological operations including opening, closing, hole-filling, and connected component filtering, yielding a binary mask $P^{\text{class}}(\mathbf{x}) \in \{0, 1\}$.

\textbf{Mask Fusion.} The final binary CoW mask $\mathcal{M}(\mathbf{x})$ was obtained as the logical union of both predictions:

\begin{equation}
    \mathcal{M}(\mathbf{x}) = \mathbb{1}\left[P^{\text{DL}}(\mathbf{x}) \geq \tau\right] \cup P^{\text{class}}(\mathbf{x})
\end{equation}

where $\tau = 0.5$ is the binarization threshold. This union strategy ensures recall is maximized, minimizing the risk of excluding valid arterial voxels from subsequent optimization.

% FIGURE 4 — Dual-Branch Segmentation:
% Side-by-side visualization of: (a) deep learning nnU-Net segmentation overlaid on the
% angiographic volume; (b) classical Frangi/Sato vesselness map and its binarized result;
% (c) final fused CoW binary mask M. A 3D rendering of the final mask colored by CoW
% segment label (ICA, MCA, ACA, PCA, BA) is recommended.
\begin{figure}[h]
    \centering
    % Insert figure here
    \caption{Dual-Branch Segmentation. Side-by-side visualization of the deep learning segmentation, classical vesselness map, and final fused CoW binary mask $\mathcal{M}$.}
    \label{fig:segmentation}
\end{figure}

\subsection*{3.5 Network Architecture: SR4DFlowNet}

\subsubsection*{3.5.1 Input Representation}

The core predictive model is an adapted SR4DFlowNet \cite{ferdian2020,elahmar2023}, a 3D fully convolutional network inspired by enhanced residual architectures for super-resolution \cite{lim2017edsr} and designed to jointly estimate high-quality velocity components and a magnitude channel from 4D Flow MRI patches. The model accepts six scalar input channels per voxel --- the three VENC-normalized velocity components $(u, v, w)$ and their corresponding magnitude-weighted counterparts $(u_m, v_m, w_m)$ --- and internally computes three auxiliary channels:

\begin{equation}
    S = \sqrt{u^2 + v^2 + w^2}, \quad
    M = \sqrt{u_m^2 + v_m^2 + w_m^2}, \quad
    \text{PCMR} = M \cdot S
\end{equation}

where $S$ is the voxel-wise flow speed, $M$ is the local signal magnitude, and PCMR (Phase-Contrast MR) captures their joint information.

\subsubsection*{3.5.2 Dual-Stream Feature Extraction}

The architecture processes velocity and physical information through two dedicated convolutional streams, each comprising two consecutive 3D convolutional layers with 64 feature maps, $3\times3\times3$ kernels, symmetric (replication) padding, and ReLU activation:

\begin{itemize}
    \item \textbf{Phase stream}: processes the velocity components $(u, v, w)$ to extract flow-kinematic features.
    \item \textbf{PC stream}: processes the physical channels $(\text{PCMR}, M, S)$ to extract signal-intensity features.
\end{itemize}

The feature maps from both streams are concatenated along the channel dimension to form a 128-channel joint representation, which is then projected to 64 channels via a $1\times1\times1$ convolution followed by a $3\times3\times3$ refinement convolution.

\subsubsection*{3.5.3 Residual Processing and Upsampling}

The fused feature map is passed through a stack of $N_{LR} = 8$ residual blocks at the input resolution \cite{he2016resnet}. Each residual block consists of two $3\times3\times3$ convolutional layers with ReLU activation and a direct identity skip connection:

\begin{equation}
    \mathbf{F}^{(l+1)} = \mathbf{F}^{(l)} + f\left(\mathbf{F}^{(l)}; \Theta^{(l)}\right)
\end{equation}

Spatial upsampling is then performed via trilinear interpolation with scale factor $r \in \{1, 2\}$ (depending on task configuration), followed by $N_{HR} = 4$ additional high-resolution residual blocks operating at the target resolution.

\subsubsection*{3.5.4 Output Heads and Magnitude Prediction}

The network concludes with three independent output heads, each comprising two convolutional layers, to predict the super-resolved velocity components $(\hat{u}, \hat{v}, \hat{w})$. When enabled, a fourth head predicts a reconstructed magnitude channel $\hat{M}$, allowing the model to simultaneously perform domain translation and velocity estimation.

\subsubsection*{3.5.5 Attention-Enhanced Variants}

Several architectural extensions incorporating attention mechanisms were evaluated. The Convolutional Block Attention Module (CBAM) \cite{woo2018cbam} applies sequential channel and spatial attention to the feature maps. Given an intermediate feature map $\mathbf{F} \in \mathbb{R}^{C \times D \times H \times W}$, CBAM computes:

\begin{equation}
    \mathbf{F}' = \mathbf{F} \otimes \mathbf{A}_c(\mathbf{F}), \qquad
    \mathbf{F}'' = \mathbf{F}' \otimes \mathbf{A}_s(\mathbf{F}')
\end{equation}

where $\mathbf{A}_c$ and $\mathbf{A}_s$ denote the channel and spatial attention gates, respectively, computed via pooled MLP and convolutional operations with sigmoid activation. More advanced variants incorporate \textbf{token-mixing attention} \cite{vaswani2017attention} --- compressing spatial features to a $P^3$ token grid via adaptive average pooling and applying multi-head self-attention in that compressed space --- and \textbf{bidirectional cross-attention} between the phase and PC feature streams, enabling explicit cross-modal feature interaction. The evaluated variants are summarized in Table~\ref{tab:variants}.

% FIGURE 5 — Network Architecture Diagram:
% Full architectural diagram of SR4DFlowNet showing: the dual-stream input (phase and PC paths),
% feature concatenation, the low-resolution residual stack, the trilinear upsampling block,
% the high-resolution residual stack, and the output heads for velocity and magnitude.
% Insets illustrating the CBAM block structure and a simplified cross-attention diagram.
\begin{figure}[h]
    \centering
    % Insert figure here
    \caption{Network Architecture Diagram of SR4DFlowNet showing the dual-stream input, residual stacks, upsampling, and output heads.}
    \label{fig:architecture}
\end{figure}

% TABLE 1 — Summary of evaluated SR4DFlowNet architectural variants and their attention components.
\begin{table}[h]
    \centering
    \caption{Summary of evaluated SR4DFlowNet architectural variants and their attention components.}
    \label{tab:variants}
    % Insert table content here
\end{table}

\subsection*{3.6 Loss Function}

Training minimized a composite loss function $\mathcal{L}$ that jointly penalizes velocity reconstruction error, magnitude reconstruction error, and texture regularity outside the vascular mask:

\begin{equation}
    \boxed{\mathcal{L} = \mathcal{L}_{\text{vel}} + \lambda_M \, \mathcal{L}_{\text{mag}} + \lambda_{\text{TV}} \, \mathcal{L}_{\text{TV}}}
\end{equation}

\subsubsection*{3.6.1 Masked Velocity Loss $\mathcal{L}_{\text{vel}}$}

The velocity loss separates contributions from intraluminal (fluid) and extraluminal (background) regions using the binary mask $\mathcal{M}$. For a predicted velocity component $\hat{q} \in \{\hat{u}, \hat{v}, \hat{w}\}$ and its normalized target $q^* = q / \text{VENC}$, the per-voxel mean squared error is:

\begin{equation}
    \epsilon_q(\mathbf{x}) = \left(\hat{q}(\mathbf{x}) - q^*(\mathbf{x})\right)^2
\end{equation}

The masked fluid MSE for each component is:

\begin{equation}
    \mathcal{L}^q_{\text{fluid}} = \frac{\sum_{\mathbf{x}} \epsilon_q(\mathbf{x}) \cdot \mathcal{M}(\mathbf{x})}{\sum_{\mathbf{x}} \mathcal{M}(\mathbf{x}) + \varepsilon}
\end{equation}

Similarly, the non-fluid MSE for each component is:

\begin{equation}
    \mathcal{L}^q_{\text{bg}} = \frac{\sum_{\mathbf{x}} \epsilon_q(\mathbf{x}) \cdot (1 - \mathcal{M}(\mathbf{x}))}{\sum_{\mathbf{x}} (1 - \mathcal{M}(\mathbf{x})) + \varepsilon}
\end{equation}

The total velocity loss aggregates over all three components with a background downweighting factor $\lambda_{\text{bg}}$:

\begin{equation}
    \mathcal{L}_{\text{vel}} = \sum_{q \in \{u,v,w\}} \left( \mathcal{L}^q_{\text{fluid}} + \lambda_{\text{bg}} \, \mathcal{L}^q_{\text{bg}} \right)
\end{equation}

where $\lambda_{\text{bg}} = 0.1$ was used throughout, reflecting the primary clinical interest in accurate intraluminal velocity quantification while still regularizing background predictions.

\subsubsection*{3.6.2 Magnitude Loss $\mathcal{L}_{\text{mag}}$}

When the magnitude head is active, an analogous masked MSE is computed on the intensity-normalized magnitude channel. The predicted magnitude $\hat{M}$ is compared against the normalized target $M^* \in [0, 1]$ obtained via per-frame min-max scaling:

\begin{equation}
    \mathcal{L}_{\text{mag}} =
        \frac{\sum_{\mathbf{x}} \left(\hat{M}(\mathbf{x}) - M^*(\mathbf{x})\right)^2 \cdot \mathcal{M}(\mathbf{x})}{\sum_{\mathbf{x}} \mathcal{M}(\mathbf{x}) + \varepsilon}
        + \lambda_{\text{bg}}
        \frac{\sum_{\mathbf{x}} \left(\hat{M}(\mathbf{x}) - M^*(\mathbf{x})\right)^2 \cdot (1-\mathcal{M}(\mathbf{x}))}{\sum_{\mathbf{x}} (1-\mathcal{M}(\mathbf{x})) + \varepsilon}
\end{equation}

The magnitude loss weight $\lambda_M = 1.0$ balanced its contribution relative to the velocity loss.

\subsubsection*{3.6.3 Outside-Mask Total Variation Penalty $\mathcal{L}_{\text{TV}}$}

To suppress spurious high-frequency texture in regions outside the vascular mask --- which would otherwise be unconstrained and produce physiologically implausible predictions --- an isotropic total variation (TV) penalty \cite{rudin1992tv,liu2019tvdip} was imposed on extraluminal voxels. For a predicted output channel $\hat{q}$, the TV penalty over the three Cartesian directions $d \in \{x, y, z\}$ is:

\begin{equation}
    \mathcal{L}_{\text{TV}} = \sum_{q} \sum_{d \in \{x,y,z\}} \frac{\sum_{\mathbf{x}} \left|\hat{q}(\mathbf{x} + \mathbf{e}_d) - \hat{q}(\mathbf{x})\right| \cdot \Omega(\mathbf{x}, d)}{\sum_{\mathbf{x}} \Omega(\mathbf{x}, d) + \varepsilon}
\end{equation}

where $\mathbf{e}_d$ is the unit displacement in direction $d$, and the joint outside indicator is:

\begin{equation}
    \Omega(\mathbf{x}, d) = \left(1 - \mathcal{M}(\mathbf{x})\right) \cdot \left(1 - \mathcal{M}(\mathbf{x} + \mathbf{e}_d)\right)
\end{equation}

The TV weight was set to $\lambda_{\text{TV}} = 10^{-5}$, providing gentle smoothness encouragement without over-regularizing the extraluminal signal.

\subsubsection*{3.6.4 Evaluation Metrics}

Model performance was assessed using two complementary metrics computed exclusively within the vascular mask. The \textbf{relative speed error} (RSE, reported as percentage) quantifies the vector-magnitude accuracy of the predicted velocity field:

\begin{equation}
    \text{RSE}(\%) = 100 \cdot \frac{
        \sum_{\mathbf{x}} \dfrac{\sqrt{(\hat{u}-u^*)^2+(\hat{v}-v^*)^2+(\hat{w}-w^*)^2}}{\sqrt{u^{*2}+v^{*2}+w^{*2}}+\varepsilon} \cdot \mathcal{M}(\mathbf{x})
    }{
        \sum_{\mathbf{x}} \mathcal{M}(\mathbf{x})
    }
\end{equation}

The \textbf{relative magnitude error} (RME) assesses the accuracy of the predicted intensity channel:

\begin{equation}
    \text{RME} = \frac{
        \sum_{\mathbf{x}} \dfrac{|\hat{M}(\mathbf{x}) - M^*(\mathbf{x})|}{|M^*(\mathbf{x})| + \varepsilon} \cdot \mathcal{M}(\mathbf{x})
    }{
        \sum_{\mathbf{x}} \mathcal{M}(\mathbf{x})
    }
\end{equation}

Both metrics were clipped to $[0, 1]$ per voxel prior to averaging to prevent outlier dominance.

% FIGURE 6 — Loss Function Schematic:
% Conceptual diagram showing: the binary mask M partitioning the volume into fluid and
% background regions; arrows indicating the three loss terms (L_vel, L_mag, L_TV) and their
% respective spatial domains; a bar chart or heatmap of per-voxel loss weights to visually
% communicate the lambda_bg = 0.1 background down-weighting.
\begin{figure}[h]
    \centering
    % Insert figure here
    \caption{Loss Function Schematic. Conceptual diagram showing the binary mask $\mathcal{M}$ partitioning the volume into fluid and background regions, and the three loss terms.}
    \label{fig:loss}
\end{figure}

\subsection*{3.7 Data Preparation and Task Configurations}

\subsubsection*{3.7.1 VENC Normalization, Magnitude Scaling, and Ground-Truth Masking}

Prior to network input, all velocity components were normalized by the global VENC:

\begin{equation}
    \tilde{q} = \frac{q}{\text{VENC}_{\text{global}}}, \quad q \in \{u, v, w\},
    \quad \text{VENC}_{\text{global}} = \max(\text{VENC}_u, \text{VENC}_v, \text{VENC}_w)
\end{equation}

yielding values approximately in $[-1, 1]$. Magnitude images were normalized to $[0, 1]$ via per-frame min-max scaling using MONAI's \texttt{ScaleIntensity} transform \cite{cardoso2022monai}, ensuring temporal consistency within each cardiac cycle.

In both task configurations described below, the reconstruction target is always the registered 7T velocity and magnitude volume \emph{after masking by} $\mathcal{M}(\mathbf{x})$:

\begin{equation}
    \mathbf{V}^{*}(\mathbf{x}, t) = \mathbf{V}_{7T}^{\text{reg}}(\mathbf{x}, t) \cdot \mathcal{M}(\mathbf{x})
\end{equation}

Applying the CoW mask to the ground truth confines supervision to the haemodynamically relevant intraluminal region, preventing the network from wasting representational capacity on extravascular tissue and ensuring that the loss is computed exclusively over voxels where accurate velocity quantification is clinically meaningful.

\subsubsection*{3.7.2 Task Configuration I --- Same-Resolution Domain Translation ($\times$1)}

In the first configuration, the network learns a direct mapping from the registered 3T acquisition to the masked 7T target at identical spatial resolution ($r = 1$). Formally, the input is the VENC-normalised 3T velocity and magnitude volume $\mathbf{V}_{3T}^{\text{reg}}(\mathbf{x}, t)$, and the target is $\mathbf{V}^{*}(\mathbf{x}, t)$. The two volumes occupy the same voxel grid by construction of the inter-scan registration pipeline (Section~3.2.2). This task constitutes a joint cross-field-strength domain translation and denoising problem: the model must simultaneously bridge the SNR gap between 3T and 7T, correct for residual contrast differences arising from distinct hardware and reconstruction chains, and predict the finer vascular detail visible only at 7T. The residual-skip variant of SR4DFlowNet was employed in this configuration, encouraging the network to learn incremental corrections relative to a direct identity shortcut rather than a full signal synthesis.

\subsubsection*{3.7.3 Task Configuration II --- Super-Resolution ($\times$2)}

In the second configuration, the network is trained to jointly perform cross-field-strength translation \emph{and} spatial super-resolution ($r = 2$): the input is a spatially downsampled version of the 3T acquisition, and the target is again the full-resolution masked 7T volume $\mathbf{V}^{*}(\mathbf{x}, t)$. The 3T volumes are downsampled by a factor of 0.5 along all three spatial dimensions using trilinear interpolation (zoom), which corresponds to the practical scenario of acquiring lower-resolution 3T data and seeking to recover both the spatial detail and the field-strength quality of 7T in a single forward pass:

\begin{equation}
    \mathbf{V}^{\text{LR}}_{3T}(\mathbf{x}, t) =
        \mathcal{Z}_{0.5}\!\left[\mathbf{V}_{3T}^{\text{reg}}(\mathbf{x}, t)\right]
\end{equation}

where $\mathcal{Z}_{0.5}$ denotes isotropic spatial downsampling by a factor of 0.5 via trilinear interpolation. The network then maps $\mathbf{V}^{\text{LR}}_{3T}$ (at half the 3T voxel size) to $\mathbf{V}^{*}$ (at full 7T resolution), combining upsampling by $r = 2$ with domain translation in a single end-to-end pass. This configuration addresses the clinically relevant setting where only low-resolution 3T data are available and the goal is to recover the spatial and haemodynamic fidelity of 7T without a second acquisition.

% FIGURE 7 — Task Configurations:
% Two-panel schematic diagram:
% (a) x1 task: 3T input (full res) --> SR4DFlowNet --> masked 7T target (full res).
%     Show example axial velocity slices and difference map between input and target.
% (b) x2 task: 3T input --> zoom x0.5 --> LR 3T input (half res) --> SR4DFlowNet (with
%     x2 upsampling) --> masked 7T target (full res).
%     Arrows should indicate the zoom step, the network upsampling, and the masking step.
%     Velocity component maps (u, v, w) and magnitude shown for both configurations.
\begin{figure}[h]
    \centering
    % Insert figure here
    \caption{Task Configurations. (a)~$\times$1 domain translation: full-resolution 3T input mapped to masked 7T target at the same resolution. (b)~$\times$2 super-resolution: trilinearly downsampled 3T input ($\times 0.5$) mapped to the masked full-resolution 7T target via a network with $r=2$ upsampling. In both cases the ground truth is the registered 7T volume restricted to the CoW mask $\mathcal{M}$.}
    \label{fig:tasks}
\end{figure}

\subsection*{3.8 Data Augmentation}

To improve generalization given the limited patient cohort, a comprehensive augmentation strategy was applied online during training \cite{zhao2019augmentation}.

\textbf{Vector-Aware Geometric Augmentation.} Random 90\textdegree{} rotations were applied with probability $p = 0.5$ per training sample. Rotation planes were selected uniformly from $\{xy, xz, yz\}$, and rotation multiplicities from $\{1, 2, 3\}$. Critically, the three velocity components were transformed as a proper vector field --- i.e., their values were permuted and sign-corrected in accordance with the rotation --- while the binary mask and magnitude images were treated as scalar fields and rotated without sign modification.

\textbf{Synthetic Nonvascular Noise Augmentation.} To prevent the network from overfitting to specific noise patterns in the background and to encourage mask-focused learning, synthetic noise was injected into the extraluminal region of the low-resolution inputs. Noise amplitudes were drawn from one of several statistical distributions (zero-mean normal, Student-$t$, Laplace, skew-normal, or generalized normal) with randomized scale parameters in each training iteration. The noise was applied exclusively outside the binary mask $\mathcal{M}$ according to:

\begin{equation}
    \tilde{\mathbf{V}}^{\text{LR}}(\mathbf{x}) = \mathbf{V}^{\text{LR}}(\mathbf{x}) + \eta(\mathbf{x}) \cdot (1 - \mathcal{M}(\mathbf{x}))
\end{equation}

where $\eta(\mathbf{x}) \sim \mathcal{D}(0, \sigma_\eta)$ with $\sigma_\eta \in [\sigma_{\min}, \sigma_{\max}]$ sampled uniformly in each batch.

\subsection*{3.9 Training Protocol}

\subsubsection*{3.9.1 Patch-Based Sampling}

All training was conducted on three-dimensional patches of size $16^3$ voxels (LR space), corresponding to $32^3$ (HR space) for the $\times$2 configuration \cite{peiffer2024patch}. During training, patches were sampled randomly from the full volumetric time series, with one cardiac frame selected uniformly at random per sample. For each training case, 64 patches were drawn per epoch; for validation, 16 center patches were used deterministically. An optional minimum mask coverage criterion ensured that sampled patches contained a sufficient proportion of vascular voxels.

\subsubsection*{3.9.2 Optimizer and Learning Rate Schedule}

The network was optimized using the Adam algorithm \cite{kingma2015adam} with an initial learning rate of $\eta_0 = 2\times10^{-4}$ and weight decay $\lambda_{L_2} = 5\times10^{-7}$. A ReduceLROnPlateau scheduler monitored the validation loss with a patience of 10 epochs, reducing the learning rate by a factor of 0.5 (minimum $\eta_{\min} = 10^{-6}$) when no improvement was observed. Early stopping was applied with a patience of 20 epochs to prevent overfitting, and an additional overfitting detector terminated training if the training loss decreased while validation loss increased for more than 15 consecutive epochs.

\subsubsection*{3.9.3 Implementation Details}

All models were implemented in PyTorch and trained on an NVIDIA GPU with a batch size of 4. The best model checkpoint (minimum validation loss) was retained for inference. Model checkpoints stored the full training state --- network weights, optimizer state, scheduler state, and hyperparameter configuration --- enabling reproducible evaluation.

% FIGURE 8 — Training Curves:
% Multi-panel figure showing training and validation loss curves over epochs for the x1 and
% x2 configurations, with learning rate adjustments annotated. A secondary panel showing the
% RSE (%) metric over epochs for both train and validation sets. Shaded regions indicating
% the overfitting detection window and early stopping point for at least one representative run.
\begin{figure}[h]
    \centering
    % Insert figure here
    \caption{Training Curves. Multi-panel figure showing training and validation loss curves for the $\times$1 and $\times$2 configurations with learning rate adjustments annotated.}
    \label{fig:training}
\end{figure}

\subsection*{3.10 Inference}

At inference time, full 4D volumes were processed using an overlap-aware sliding-window strategy implemented via MONAI's \texttt{sliding\_window\_inference} \cite{cardoso2022monai,ronneberger2015unet}. For each time frame, the volume was partitioned into overlapping $16^3$ patches (25\% overlap, i.e., stride of $12$ voxels per dimension), and each patch was independently processed by the frozen network. Predicted patches were recombined using Gaussian-weighted blending at overlapping boundaries to eliminate patch-edge discontinuities:

\begin{equation}
    \hat{\mathbf{V}}(\mathbf{x}) = \frac{\sum_k w_k(\mathbf{x}) \cdot \hat{\mathbf{V}}_k(\mathbf{x})}{\sum_k w_k(\mathbf{x})}
\end{equation}

where $w_k(\mathbf{x})$ is a Gaussian window centered on patch $k$, and the sum is over all patches covering voxel $\mathbf{x}$. The resulting full-volume predictions were post-masked by $\mathcal{M}(\mathbf{x})$ to constrain the output to the CoW region, and the predicted velocity components were rescaled to physical units via $\hat{q}_{\text{phys}} = \hat{q} \cdot \text{VENC}$.

% FIGURE 9 — Inference Pipeline and Qualitative Results:
% Comprehensive results figure showing: (a) velocity vector maps (mid-systolic frame) from the
% 3T input, the 7T reference, and the network prediction for x1 and x2 configurations on
% representative axial slices through the Circle of Willis; (b) voxel-wise relative speed error
% maps overlaid on the angiographic volume; (c) 3D streamline or arrow plots of the reconstructed
% velocity field within the CoW mask. Zoomed inset on the MCA bifurcation recommended.
\begin{figure}[h]
    \centering
    % Insert figure here
    \caption{Inference Pipeline and Qualitative Results. Velocity vector maps, speed error maps, and 3D streamline plots for the $\times$1 and $\times$2 configurations.}
    \label{fig:inference}
\end{figure}

% ── Bibliography ──────────────────────────────────────────────────────────────
% Use the following in your main .tex file:
%   \bibliographystyle{ieeetr}   % or unsrt, apalike, etc.
%   \bibliography{neuroflow}
%
% All entries are defined in neuroflow.bib (same directory).
\bibliographystyle{ieeetr}
\bibliography{neuroflow}
